% Generated by task tex-groundtruth from dossiers/gt-table-pub.md (sections 1, 2 and 7),
% with the corrections of dossiers/verify-gt-table-pub.md.
\section{The table sumcheck and the public input: the transcript as deployed}\label{sec:gen-table-pub-transcript}

Short forms used in this section. The number after the colon is the line, at leanVM's pin
\code{a386121f}; a bare \code{:N} is a line of the file last named.

{\small\begin{tabular}{@{}ll@{}}
\code{03:}, \code{05:}, \code{07:}, \code{08:} & \file{doc/leanvm/body/03-proving-primitives.tex}, \file{05-arithmetization.tex}, \\
 & \file{07-instruction-tables.tex}, \file{08-end-to-end-protocol.tex} \\
Rust files & under \file{crates/lean_vm/src/}, except the following \\
\code{transcript.rs} & \file{crates/fiat_shamir/src/transcript.rs} \\
\code{stack_open.rs} & \file{crates/pcs/src/stack_open.rs} \\
\code{primitives/src/…} & \file{crates/primitives/src/…} \\
\code{py:} & \file{python-verifier/verifier.py} \\
\code{aggregate.py} & \file{crates/rec_aggregation/guests/aggregate.py}, the recursion guest, \\
 & a third executable verifier \\
\end{tabular}}

\subsection{Notation}\label{sec:gen-table-pub-notation}

$\K = \mathrm{GF}(2^{64})$, $\E = \mathrm{GF}(2^{192}) = \K[y]/(y³+y+1)$. Six opcode tables
$t = 0..5$ in the order XOR, MUL\_NATIVE, SET\_CONSTANT, DEREF, JUMP, BLAKE2S, of log-heights
$τ_t$ and widths $n_t = 15, 15, 8, 15, 14, 37$ (sum \textbf{104}). $τ_{\max} = \max_t τ_t$ (at
least 3: the BLAKE2S table has $τ_5 ≥ 3$). $B = Σ_t (\text{number of constraints of } t) = 2$
(JUMP's two; every other table has none). Three \emph{sides} $s = 0, 1, 2$: push, pull, count.
$ζ ∈ \E^{μ_{\mathrm{bus}}}$ is the terminal point of the bus phase's GKR. $\mathit{totals}_s$
(the specification's $\mathit{rem}_s$) is what the tables owe side $s$. $B^s_t$ is table $t$'s
bus form for side $s$, a polynomial of degree at most 2 in the columns of $t$ with coefficients
in $\E$.

Widths, from \code{tables.rs}: \code{arith::N = 15} (\code{:455}), \code{set::N = 8}
(\code{:537}), \code{deref::N = 15} (\code{:608}), \code{jump::N = 14} (\code{:711}),
\code{blake2st::N = 37} (\code{:862}); Python \code{TABLE_WIDTHS} (\code{py:823-834, 852});
specification §7 column lists (\code{07:12, 30, 58, 78, 99, 122}), which give the same counts
(BLAKE2S: 19 stacked columns plus the 18 limbs).

\subsection{The table sumcheck}\label{sec:gen-table-pub-transcript-table}

\begingroup\footnotesize\setlength{\tabcolsep}{4pt}\renewcommand{\arraystretch}{1.15}
\begin{longtable}{@{}>{\raggedright\arraybackslash}p{0.04\linewidth}>{\raggedright\arraybackslash}p{0.295\linewidth}>{\raggedright\arraybackslash}p{0.19\linewidth}>{\raggedright\arraybackslash}p{0.24\linewidth}>{\raggedright\arraybackslash}p{0.16\linewidth}@{}}
\caption{The table sumcheck, step by step, in the three sources.}\label{tab:gen-table-pub-transcript-table}\\
\toprule
\# & Step & Specification & Rust & Python \\
\midrule
\endfirsthead
\toprule
\# & Step & Specification & Rust & Python \\
\midrule
\endhead
\bottomrule
\endfoot
T0 & What precedes: the bus phase's last prover message is the five boundary evaluations
  (\code{mem₀}, \code{mem₁}, \code{mem₂}, \code{cntfin_mem} at $ζ_{<κ_{\mathrm{mem}}}$;
  \code{cntfin_bc} at $ζ_{<κ_{\mathrm{bc}}}$)
  & \code{08:70} ``For their committed columns the prover sends one evaluation each''
  & \code{leaf.rs:909-920} (\code{decompose_verify}), order fixed by \code{cpu/layout.rs:361-395}
  & \code{py:552-557} \\
T1 & \textbf{Challenge} $ξ ← \E$ (one element), drawn after the boundary evaluations (T0)
  & \code{08:75} ``Samples and sends the batch challenge $ξ∈\E$''; \code{05:145}
  & \code{cpu/mod.rs:726} \code{let zc_xi = vs.sample();} (prover \code{:590})
  & \code{py:1389} \\
T2 & \textbf{Powers.} Constraint $i$ (0-based) of table $t$ takes $ξ^{o_t+i}$, $o_t$ the number of
  constraints of the tables before $t$; side $s$ takes $ξ^{B+s}$, shared by every table. For
  leanVM: $ξ⁰ → b + v_{\mathrm{cond}}·w$, $ξ¹ → v_{\mathrm{cond}}·(b+1)$, $ξ², ξ³, ξ⁴ →$ push,
  pull, count
  & \code{05:145}; \code{08:75} ``a disjoint range per table; the last $ν_{\mathrm{side}}$ weight
  the bus forms, one per side and shared by every table''
  & \code{constraints.rs:74-82} (\code{xi_offsets}), \code{:254-255, 281};
  \code{cpu/mod.rs:413-422} (\code{xi_form_base}, \code{xi_form_pows}); \code{tables.rs:70-74}
  & \code{py:1390-1392}, \code{:621-623}, \code{:798-800} \\
T3 & \textbf{Target, derived, never sent}: $T = Σ_s ξ^{B+s}·\mathit{totals}_s$
  & \code{05:151-155} ``the verifier computing the right side itself''; \code{08:76} ``against
  the target the verifier derived from the three leaf claims''
  & \code{cpu/mod.rs:735}
  & \code{py:1393} \\
T4 & \textbf{Number of rounds} $τ_{\max}$, the maximum over the six opcode tables
  & \code{05:145, 155}; \code{08:76}
  & \code{constraints.rs:250} \code{airs.iter().map(|a| a.tau).max()}
  & \code{py:607} \code{max(table_log_heights)} \\
T5 & \textbf{Round $j = 0..τ_{\max}−1$ binds variable $m = τ_{\max}−1−j$} (highest first).
  \textbf{Prover message}: three elements of $\E$, the coefficients $c₀, c₂, c₃$ of the cubic
  $h_j$, in that order
  & \code{05:155} ``Rounds bind $X_{τ_{\max}−1}$ first''; \code{08:76} ``the round check pins
  one of the four coefficients (§3), so a round costs three $\E$-elements'' (which coefficient:
  not said)
  & prover \code{constraints.rs:164, 191-194} (\code{add_round_poly(&h, false)}); encoding
  \code{transcript.rs:63-71} (\code{fixed = 1})
  & \code{py:409} ``The transcript contains c0, c2, c3'' \\
T6 & \textbf{Verifier computation}: $c₁ := \mathit{claim} + c₂ + c₃$, so that
  $h_j(0)+h_j(1) = \mathit{claim}$ holds by construction. No rejection path
  & \code{03:84} states it as a check, \code{03:110} as the reason one coefficient is not sent
  & \code{transcript.rs:289-309}; \code{constraints.rs:267}
  & \code{py:409-410} \\
T7 & \textbf{Challenge} $r_m ← \E$ after the round message; $\mathit{claim} := h_j(r_m)$
  & \code{03:84}
  & \code{constraints.rs:268-270}
  & \code{py:424-426} \\
T8 & \textbf{Joining and padding.} Table $t$ is active in the round binding $m$ iff $τ_t > m$,
  that is from round $τ_{\max} − τ_t$ on. Verifier's weight:
  $w_t := ∏_{m<τ_t}(1+ζ_m+r_m) · ∏_{τ_t ≤ m < τ_{\max}} r_m$
  & \code{05:146-150, 155}
  & \code{constraints.rs:271-274}
  & \code{py:609-614} \\
T9 & \textbf{Final prover message}: for $t = 0..5$ in order, $n_t$ elements: the value of every
  column of $t$, the eighteen BLAKE2S limbs included, at $r_{<τ_t} = (r_0,…,r_{τ_t−1})$. 104
  elements, no challenge in between
  & \code{08:77} ``The prover sends one value per column of every table''
  & \code{constraints.rs:224-238} (prover), \code{:279-280} (verifier); \code{cpu/mod.rs:378}
  \code{n_cols: table.n_committed_columns()}
  & \code{py:619-620} \\
T10 & \textbf{Final check}:
  $Σ_t w_t·(Σ_i ξ^{o_t+i}·C_{t,i}(e_t) + Σ_s ξ^{B+s}·B^s_t(e_t)) = \mathit{claim}$, else reject
  & \code{08:77} ``checks that they reproduce the sumcheck's final value''
  & \code{constraints.rs:277-290} (\code{Error::FinalMismatch}); summand
  \code{cpu/mod.rs:380-383}
  & \code{py:616-625} \\
T11 & \textbf{Claims pooled}: for $t$, then for local column $c$: column $\mathit{base}_t + c$,
  point $r_{<τ_t}$, value $e_{t,c}$
  & \code{08:77} ``the column values then enter the claim pool, $T_j$'s at $ρ_{<τ_j}$''
  & \code{cpu/mod.rs:428-441} (\code{constraint_claims})
  & \code{py:624} \\
T12 & \textbf{The eighteen limbs}: their claims are in the pool at their place (BLAKE2S local
  columns 9 to 26) and are \emph{located} on $\qflock$: low 8 coordinates frozen to the slot's
  bits, then the point, then $\qflock$'s selector
  & \code{08:77} ``live in $\qflock$, not the stack, so their claims route there'';
  \code{07:120-122}
  & \code{cpu/mod.rs:447-453, 790-814} (\code{SlotClaim::Strided}); \code{tables.rs:394-413};
  \code{hash_flock.rs:96-115}
  & \code{py:847-850, 886-894}; \code{:295-297} \\
\end{longtable}
\endgroup

\begin{table}[h]
\centering\small
\begin{tabular}{@{}lll@{}}
\toprule
Phase & Elements of $\E$ sent by the prover & Challenges drawn by the verifier \\
\midrule
table sumcheck & $3·τ_{\max} + 104$ & $1 + τ_{\max}$ \\
public input & 2 (step P2) & 1 (step P1) \\
\bottomrule
\end{tabular}
\caption{Counts of the two phases.}\label{tab:gen-table-pub-counts}
\end{table}

Verbatim, the Rust verifier's loop and final check (\code{constraints.rs:261-291}):

\begin{srccode}
    let mut claim = target;
    let mut chi = vec![F192::ZERO; n];
    for j in 0..n {
        let m = n - 1 - j;
        // `h(0)` is derived from the running claim rather than transmitted, so
        // the round-consistency check it used to enable holds by construction.
        let h = vs.next_round_poly(4, claim, None).map_err(|_| Error::Truncated)?;
        let rk = vs.sample();
        chi[m] = rk;
        claim = poly_eval(&h, rk);
        let eq_k = F192::ONE + zeta[m] + rk;
        for (t, air) in airs.iter().enumerate() {
            weights[t] *= if air.tau > m { eq_k } else { rk };
        }
    }

    let mut acc = F192::ZERO;
    let mut claims = Vec::with_capacity(airs.len());
    for (t, air) in airs.iter().enumerate() {
        let evals = vs.next_scalars(air.n_cols).map_err(|_| Error::Truncated)?;
        let w = &pows[offsets[t]..offsets[t] + air.n_constraints];
        acc += weights[t] * (air.eval)(w, &evals);
        claims.push(Claims {
            chi: chi[..air.tau].to_vec(),
            evals,
        });
    }
    if acc != claim {
        return Err(Error::FinalMismatch);
    }
    Ok(claims)
\end{srccode}
\src{crates/lean\_vm/src/constraints.rs:261-291}

and the decoding (\code{transcript.rs:289-302}): with \code{eq = None}, \code{fixed = 1}, the
three transmitted coefficients are read into indices 0, 2, 3 and
\code{coeffs[fixed] = claim + sum_from(2)}.

\subsection{The public input}\label{sec:gen-table-pub-transcript-pub}

\begingroup\footnotesize\setlength{\tabcolsep}{4pt}\renewcommand{\arraystretch}{1.15}
\begin{longtable}{@{}>{\raggedright\arraybackslash}p{0.04\linewidth}>{\raggedright\arraybackslash}p{0.225\linewidth}>{\raggedright\arraybackslash}p{0.22\linewidth}>{\raggedright\arraybackslash}p{0.27\linewidth}>{\raggedright\arraybackslash}p{0.17\linewidth}@{}}
\caption{The public-input phase, step by step, in the three sources.}\label{tab:gen-table-pub-transcript-pub}\\
\toprule
\# & Step & Specification & Rust & Python \\
\midrule
\endfirsthead
\toprule
\# & Step & Specification & Rust & Python \\
\midrule
\endhead
\bottomrule
\endfoot
P0 & Precondition: each public word has top limb 0
  & \code{08:29} ``two 192-bit words (with top limb 0)''; no check stated
  & type \code{[F192; 2]}; \code{read_public} rejects \code{half.c2 != 0} with
  \code{CpuError::PublicInput} (\code{cpu/mod.rs:141-143})
  & by type: \code{Digest} is 32 bytes, \code{halves()} builds \code{E(w0, w1)}
  (\code{py:204-219}) \\
P1 & \textbf{Challenge} $r ← \E$, after the 104 values of the final message (T9)
  & \code{08:83}
  & \code{cpu/mod.rs:745} (prover \code{:610})
  & \code{py:1398} \\
P2 & \textbf{Prover message}: two elements $c₀, c₁$
  & \code{08:29, 84}
  & \code{cpu/mod.rs:746-749} (prover \code{:611-618})
  & \code{py:1399} \\
P3 & \textbf{Check}
  & per limb: $c_ℓ = (1+r)·\mathrm{mem}[g⁰]_ℓ + r·\mathrm{mem}[g¹]_ℓ$, $ℓ ∈ \{0,1\}$
  (\code{08:29-32})
  & combined: $c₀ + y·c₁ = \mathrm{interp}(\mathrm{pi}₀, \mathrm{pi}₁, r)$
  (\code{cpu/mod.rs:752-755})
  & combined, the same (\code{py:1400}) \\
P4 & \textbf{Claims pooled}: \code{mem₀}, \code{mem₁}, \code{mem₂} at $(r, 0, …, 0)$ with values
  $c₀, c₁, 0$; no scalar for the third
  & \code{08:84} ``pools them, with 0 for the top limb''
  & \code{cpu/mod.rs:746} (array initialised to zero, two entries read), \code{:674-682}
  (\code{bind_pi_claim})
  & \code{py:1399, 1401-1402} \\
\end{longtable}
\endgroup

\subsection{The claim pool, in order}\label{sec:gen-table-pub-transcript-pool}

\begingroup\small\setlength{\tabcolsep}{4pt}\renewcommand{\arraystretch}{1.15}
\begin{longtable}{@{}>{\raggedright\arraybackslash}p{0.18\linewidth}>{\raggedright\arraybackslash}p{0.39\linewidth}>{\raggedright\arraybackslash}p{0.08\linewidth}>{\raggedright\arraybackslash}p{0.29\linewidth}@{}}
\caption{The order of the claim pool.}\label{tab:gen-table-pub-pool}\\
\toprule
Rank & Claims & Number & Source \\
\midrule
\endfirsthead
\toprule
Rank & Claims & Number & Source \\
\midrule
\endhead
\bottomrule
\endfoot
first power of $λ$ & the ring-switched claim on $\qflock$ & 1
  & \code{08:98}; \code{stack_open.rs:518-519}; \code{py:1413} \\
then & bus: \code{mem₀}, \code{mem₁}, \code{mem₂}, \code{cntfin_mem} at
  $ζ_{<κ_{\mathrm{mem}}}$, \code{cntfin_bc} at $ζ_{<κ_{\mathrm{bc}}}$ & 5
  & \code{cpu/mod.rs:663}; \code{py:555-557, 1395} \\
then & table sumcheck: the column claims (T11), of which 18 strided on $\qflock$ & 104
  & \code{cpu/mod.rs:664}; \code{py:1395} \\
then & public input: the pooled claims (P4) & 3 & \code{cpu/mod.rs:665}; \code{py:1402} \\
\end{longtable}
\endgroup

\code{finish_claims} (\code{cpu/mod.rs:656-667}) is the one place the order is fixed in the Rust.
The pool has 112 point claims and the batch 113 powers of $λ$.

\subsection{The public-input check, verbatim in the four sources}\label{sec:gen-table-pub-transcript-pubcheck}

\textbf{Specification} (\code{08:29-33}):

\begin{srccode}
The public input is the first two memory cells $\mem[\gen^{0}],\mem[\gen^{1}]$, two $192$-bit words (with top limb $0$) both parties know. The verifier samples $r_m\in\E$, the prover sends $c_0,c_1$ claiming $c_\ell=\mle{\mem_\ell}(r_m,0,\dots,0)$, and the verifier checks
\[
  c_\ell\;\qeq\;(1+r_m)\,\mem[\gen^{0}]_\ell+r_m\,\mem[\gen^{1}]_\ell,\qquad \ell\in\{0,1\}.
\]
Both sides are degree one in $r_m$, so a limb disagreeing with the public words survives with probability at most $1/|\E|$ (Lemma~\ref{lem:sz}). The claims $c_0$, $c_1$ are discharged to the PCS.
\end{srccode}
\src{doc/leanvm/body/08-end-to-end-protocol.tex:29-33}

\textbf{Rust} (\code{cpu/mod.rs:745-756}):

\begin{srccode}
    let r_pi = vs.sample();
    let mut pi_limbs = [F192::ZERO; 3];
    for v in &mut pi_limbs[..2] {
        *v = vs.next_scalar().map_err(CpuError::Transcript)?;
    }
    // The two claimed evaluations must sit on the public-input line, the top
    // limb's being zero (§sec:e2e-pi).
    let want = primitives::multilinear::interp(l.pi[0], l.pi[1], r_pi);
    if pi_limbs[0] + F192::Y * pi_limbs[1] != want {
        return Err(CpuError::PublicInput);
    }
    let slots = finish_claims(&l, bus.claims, &table_claims, r_pi, pi_limbs);
\end{srccode}
\src{crates/lean\_vm/src/cpu/mod.rs:745-756}

with \code{interp(lo, hi, t) = lo + t * (lo + hi)} (\code{primitives/src/multilinear.rs:39-41})
and \code{Y = (0, 1, 0)} (\code{primitives/src/field/gf2_64x3.rs:44}).

\textbf{Python} (\code{py:1398-1402}):

\begin{srccode}
    public_challenge = transcript.sample()
    public_limbs = (*transcript.next_scalars(2), ZERO)
    require(poly_eval(public_limbs, Y) == multilinear_eval(public_input.halves(), [public_challenge]), "public input check failed")
    public_point = (public_challenge, *[ZERO] * (layout.placements[MEMORY_0].variables - 1))
    claims.extend(ColumnClaim(column, public_point, value) for column, value in zip((MEMORY_0, MEMORY_1, MEMORY_2), public_limbs))
\end{srccode}
\src{python-verifier/verifier.py:1398-1402}

\textbf{Recursion guest} (\code{aggregate.py:1672-1688}):

\begin{srccode}
    # ---- public-input binding claim: MEM as ONE logical E-column ----
    # The VM's bind_pi_claim makes a SINGLE E-claim at [rm, 0..]:
    #   MEM(rm) = interp(pi_0, pi_1, rm) = pi_0 + rm*(pi_0 + pi_1)
    # over the E-valued public input (no lane splitting, no Frobenius). Both
    # public words have a zero top limb, so that limb's evaluation is zero at
    # every rm; only the two low ones ride the stream and must reassemble it:
    # MEM = v_lo + Y*v_hi (doc sec:e2e-pi).
    fs, rm = squeeze(fs)
    mem = pi_0 + rm * (pi_0 + pi_1)
    fs, mem_lo, cursor = fs_next(fs, cursor)
    fs, mem_hi, cursor = fs_next(fs, cursor)
    assert mem == mem_lo + mem_hi * Y_TOWER
    claim_pool[GEN ** claim_idx] = mem_lo
    claim_idx += 1
    claim_pool[GEN ** claim_idx] = mem_hi
    claim_idx += 1
    claim_pool[GEN ** claim_idx] = 0
\end{srccode}
\src{crates/rec\_aggregation/guests/aggregate.py:1672-1688}

\paragraph{Which transcripts each accepts.} Write $w_i = w_{i,0} + y·w_{i,1}$ for the public words
and $L_ℓ(r) = (1+r)·w_{0,ℓ} + r·w_{1,ℓ}$, elements of $\E$. Then
$\mathrm{interp}(w_0, w_1, r) = L_0(r) + y·L_1(r)$.

\begin{table}[h]
\centering\small
\begin{tabular}{@{}>{\raggedright\arraybackslash}p{0.2\linewidth}>{\raggedright\arraybackslash}p{0.32\linewidth}>{\raggedright\arraybackslash}p{0.4\linewidth}@{}}
\toprule
 & Accepts $(r, c₀, c₁)$ iff & Pairs per $r$ \\
\midrule
Specification & $c₀ = L_0(r)$ and $c₁ = L_1(r)$ & 1 \\
Rust, Python, guest & $c₀ + y·c₁ = L_0(r) + y·L_1(r)$ & $|\E| = 2^{192}$: the pairs
  $(L_0(r) + y·δ, L_1(r) + δ)$, $δ ∈ \E$ \\
\bottomrule
\end{tabular}
\caption{The transcripts the per-limb check and the combined check accept.}\label{tab:gen-table-pub-accepts}
\end{table}

The first implies the second. Not conversely: $c₀, c₁$ are elements of $\E$, so the independence
of $1, y$ over $\K$ says nothing of them.

\paragraph{With the pooled claims true of a committed memory.} Let the cells differ from the
public words by $a_ℓ$ (cell 0) and $c_ℓ$ (cell 1), in $\K$. The true evaluations pass
\begin{itemize}
\item the per-limb check iff $(1+r)·a_ℓ + r·c_ℓ = 0$ for $ℓ = 0$ and $ℓ = 1$;
\item the combined check iff $(1+r)·A + r·C = 0$, $A = a_0 + y·a_1$, $C = c_0 + y·c_1$.
\end{itemize}
If the memory is wrong, $(A, C) ≠ (0, 0)$ because $1, y$ are independent over $\K$, and
$A + r·(A + C) = 0$ has at most one solution ($A + C = 0$ forces $A = 0$, then $C = 0$). So the
combined check has knowledge error $1/|\E|$ too. The bad challenges differ: with
$a_0 = 1, c_0 = 0, a_1 = 1, c_1 = 1 + g$ the per-limb check has none (limb 0 needs $r = 1$,
limb 1 needs $r = g^{-1}$) and the combined check has $r = (1+y)/(1+y·g)$. The status's example
(zero words, limbs $[0,1]$ and $[1,0]$, $r = y/(1+y)$) is checked by
\file{tests/LeanerVMTests/Protocol/PublicInput.lean:182-210} (leanerVM \code{b435631}); the
dossier's author re-derived it by hand.

\paragraph{Is the Rust's public input two words of $\E$ or of 128 bits?} By type two words of
$\E$ (\code{[F192; 2]}, \code{Layout.pi}, \code{cpu/layout.rs:83}). By what the verifier accepts,
two words of 128 bits: \code{read_public} rejects a nonzero top limb, and the transcript is
seeded with four lanes.

\paragraph{Which is the deployed protocol.} The combined check. Three executable verifiers at the
pin run it, the recursion guest among them, and the guest's comment states the intent (``MEM as
ONE logical E-column''). The specification's §8.2 is the outlier. An honest Rust proof passes
both: the Rust prover sends $L_0(r), L_1(r)$ (\code{cpu/mod.rs:611-613}).
